\begin{tikzpicture}[line width=.6pt, font=\normalsize]
% cone body: top rim ellipse rx=3, apex below
\coordinate (rimL) at (-3,0);
\coordinate (rimR) at (3,0);
\coordinate (apex) at (0,-6);
\coordinate (wL) at (-1.5,-3);
\coordinate (wR) at (1.5,-3);
% water region (gray block with boundary and no hole issue)
\fill[gray!15] (wL) -- (apex) -- (wR) -- cycle;
\fill[gray!15] (wL) arc [x radius=1.5, y radius=0.28, start angle=180, end angle=360] -- cycle;
% cone outline
\draw (rimL) -- (apex) -- (rimR);
\draw (rimL) arc [x radius=3, y radius=0.55, start angle=180, end angle=360];
\draw (rimL) arc [x radius=3, y radius=0.55, start angle=180, end angle=0];
% water surface: solid front arc, dashed back arc
\draw (wL) arc [x radius=1.5, y radius=0.28, start angle=180, end angle=360];
\draw[dashed] (wL) arc [x radius=1.5, y radius=0.28, start angle=180, end angle=0];
% water label inside region
\node at (0,-3.9) {\begin{tabular}{c}雨\\水\end{tabular}};
% dimensions: top diameter 200mm
\draw[dimension line] (-3,1.1) -- node[dimension label] {200 mm} (3,1.1);
\draw[dimension extension] (-3,0.55) -- (-3,1.1);
\draw[dimension extension] (3,0.55) -- (3,1.1);
% left total height 300mm
\draw[dimension line] (-4.4,0) -- node[dimension label vertical] {300 mm} (-4.4,-6);
\draw[dimension extension] (-3,0) -- (-4.4,0);
\draw[dimension extension] (0,-6) -- (-4.4,-6);
% right water height 150mm
\draw[dimension line] (4.0,-3) -- node[dimension label vertical] {150 mm} (4.0,-6);
\draw[dimension extension] (1.5,-3) -- (4.0,-3);
\draw[dimension extension] (0,-6) -- (4.0,-6);
\end{tikzpicture}
